\chapter{2022年「数理手法 • 物理学専攻」期末试题解答}
\section{填空题（每题4分，共20分）}
{
\vspace{-1em}
\begin{exercise}
    $\left(1+\ii\right)^3$的三角表示法为\underline{$2\sqrt{2}\left(\cos{\frac{3\pi}{4}+\ii\sin{\frac{3\pi}{4}}}\right)$}.
\end{exercise}
\begin{solution}
    $1+\ii$的三角表示法为$\sqrt{2}\left(\cos{\frac{\pi}{4}+\ii\sin{\frac{\pi}{4}}}\right)$，由棣莫弗公式
    $$\left[r\left(\cos{\theta}+\ii\sin{\theta}\right)\right]^n=r^n\left(\cos{n\theta}+\ii\sin{n\theta}\right)$$
    得$\left(1+\ii\right)^3=2\sqrt{2}\left(\cos{\frac{3\pi}{4}+\ii\sin{\frac{3\pi}{4}}}\right)$.
\end{solution}
\begin{exercise}
    $\ee^z+1=0$的解为\underline{$z=\ii\left(2k+1\right)\pi$}.
\end{exercise}
\begin{solution}
    由Euler公式
    $$
    \begin{aligned}
    \ee^z&=\ee^{\ii\left(2k+1\right)\pi}\\
    &=\cos{\left(2k+1\right)\pi}+\ii\sin{\left(2k+1\right)\pi}\\
    &=-1
    \end{aligned}
    $$
    得$z=\ii\left(2k+1\right)\pi$.
\end{solution}

\begin{exercise}
    规定在$z=1$时，$\ln{z}=0$.则对复变函数$w\left(z\right)=\ee^{z\ln{z}}$，当$z$从$z=1$处发，沿下半平面到$z=-e$处，则$w\left(z\right)$在$z=-e$处的值为\underline{}.
\end{exercise}
\begin{solution}
    
\end{solution}

\begin{exercise}
    $m$阶Bessel函数的表达式为\underline{$\mathbf{J}_m\left(x\right)=\sum_{k=0}^{\infty}{\frac{\left(-1\right)^k}{k!\Gamma\left(k+\nu+1\right)}\left(\frac{x}{2}\right)^{2k+\nu}}$}.
\end{exercise}

\begin{exercise}
    $\sin{t}$的Laplace变换为\underline{$\frac{1}{1+p^2}$}.
\end{exercise}
\begin{solution}
$\sin{\omega t}$的Laplace变换为
$$
\begin{aligned}
    \mathcal{L}\left(p\right)&=\int_0^{\infty}{\sin{\omega t}\ee^{-pt}\dd t}\\
    &=-\frac{1}{p}\int_0^{\infty}{\sin{\omega t\dd\ee^{-pt}}}\\
    &=-\frac{1}{p}\left[\sin{\omega t}\ee^{-pt}\Bigg|_0^{\infty}-\int_0^{\infty}{\ee^{-pt}\omega\cos{\omega t}\dd t}\right]\\
    &=-\frac{1}{p}\left\{0-\omega\left[-\int_0^{\infty}{\frac{1}{p}\cos{\omega t}\dd\ee^{-pt}}\right]\right\}\\
    &=-\frac{\omega}{p^2}\int_0^{\infty}{\cos{\omega t}\dd\ee^{-pt}}\\
    &=-\frac{\omega}{p^2}\left[\cos{\omega t}\ee^{-pt}\Bigg|_0^{\infty}-\left(-\omega\int_0^{\infty}{\sin{\omega t}\ee^{-pt}\dd t}\right)\right]\\
    &=-\frac{\omega}{p^2}\left[0-1+\omega \mathcal{L}\left(p\right)\right]\\
    &=\frac{\omega}{p^2}-\frac{\omega^2}{p^2}\mathcal{L}\left(p\right)
\end{aligned}
$$
即$\mathcal{L}\left(p\right)=\frac{\omega^2}{\omega^2+p^2}$；令$\omega=1$得$\sin{t}$的Laplace变换为$\frac{1}{1+p^2}$.

\end{solution}
}
\section{计算题（每题6分，共30分）}
{
\begin{exercise}
    已知某复变函数的实部$u\left(x,y\right)=\ee^x\left(x\cos{y}-y\sin{y}\right)$，且$f\left(0\right)=0$.求$f\left(z\right)$的表达式.
\end{exercise}
\begin{solution}
    $$u_x=\ee^x\left(x\cos{y}-y\sin{y}+\cos{y}\right)$$
    由题意可知，$f\left(z\right)$是解析的，故满足Cauchy Riemann方程$u_x=v_y$，即
    $$\frac{\partial v}{\partial y}=\ee^x\left(x\cos{y}-y\sin{y}+\cos{y}\right)$$
    积分，得
    $$
    \begin{aligned}
    v&=\int_0^y{\ee^x\left(x\cos{y}-y\sin{y}+\cos{y}\right)}\dd y\\
    &=\ee^x\left(x\sin{y}+y\cos{y}\right)+C
    \end{aligned}
    $$
    代入条件$f\left(0\right)=0$，得$C=0$.故$f\left(z\right)$的表达式为
    $$f\left(z\right)=\ee^x\left(x\cos{y}-y\sin{y}\right)+\ii\ee^x\left(x\sin{y}+y\cos{y}\right)$$
    当然，本题求出$u_y$后利用条件$u_y=-v_x$也可以.
\end{solution}

\begin{exercise}
    写出二阶Legendre多项式的表达式和$l=2,\ m=1$时的连带Legendre函数.
\end{exercise}
\begin{solution}
    $$
    \begin{aligned}
    P_l\left(x\right)&=\frac{1}{2^ll!}\frac{\dd^l}{\dd x^l}\left(x^2-1\right)^l\\
    &\xlongequal{l=2}\frac{1}{2}\left(3x^2-1\right)
    \end{aligned}
    $$
    $$
    \begin{aligned}
    P_l^m&=\left(1-x^2\right)^{\frac{m}{2}}P_l^{\left(m\right)}\left(x\right)\\
    &\xlongequal{l=2,\ m=1}3x\left(1-x^2\right)^{\frac{1}{2}}
\end{aligned}
    $$
\end{solution}

\begin{exercise}
    求解下列PDE
    $$
    \left\{
        \begin{aligned}
            &u_{tt}=u_{xx}\\
            &u\left(x,0\right)=x^2\\
            &u_t\left(x,0\right)=x
        \end{aligned}
    \right.
    $$
\end{exercise}
\begin{solution}
    $$a=1,\ \varphi\left(x\right)=x^2,\ \psi\left(x\right)=x$$
由d'Alembert行波解
$$
\begin{aligned}
    u\left(x,t\right)&=\frac{1}{2}\left[\varphi\left(x+at\right)+\varphi\left(x-at\right)\right]+\frac{1}{2a}\int_{x-at}^{x+at}{\psi\left(\alpha\right)\dd\alpha}\\
    &=\frac{1}{2}\left[\left(x+t\right)^2+\left(x-t\right)^2\right]+\frac{1}{2}\int_{x-t}^{x+t}{\alpha\dd\alpha}\\
    &=x^2+xt+t^2
\end{aligned}
$$
\end{solution}

\begin{exercise}
    求积分$\int_0^{\infty}{\frac{1}{\left(z^2+1\right)\left(z^2+4\right)}\dd z}$.
\end{exercise}
\begin{solution}
    被积函数$f\left(z\right)=\frac{1}{\left(z^2+1\right)\left(z^2+4\right)}$存在$4$个单极点：$z_1=-\ii,\ z_2=\ii,\ z_3=-2\ii,\ z_4=2\ii$.四个极点的留数分别为
    $$\Res[f\left(z\right),z_i]=\lim_{z\to z_i}\left(z-z_i\right)f\left(z\right)$$
    即
    $$
    \begin{aligned}
        &\Res[f\left(z\right),z_i]=\lim_{z\to z_1}\left(z-z_1\right)f\left(z\right)=-\frac{1}{6\ii}&&\Res[f\left(z\right),z_i]=\lim_{z\to z_2}\left(z-z_2\right)f\left(z\right)=\frac{1}{6\ii}\\
        &\Res[f\left(z\right),z_i]=\lim_{z\to z_3}\left(z-z_3\right)f\left(z\right)=\frac{1}{12\ii}&&\Res[f\left(z\right),z_i]=\lim_{z\to z_4}\left(z-z_4\right)f\left(z\right)=-\frac{1}{12\ii}
    \end{aligned}
    $$
    由留数定理，原积分的值为
    $$I=2\pi\ii\sum_{i=1}^4{\Res[f\left(z\right),z_i]}=0$$
\end{solution}

\begin{exercise}
    计算实变函数定积分$\int_0^{2\pi}{\frac{1}{2+\cos{\theta}}\dd\theta}$.
\end{exercise}
\begin{solution}
    作变换$z=\ee^{\ii\theta}$，即$\cos{z}=\frac{z+z^{-1}}{2},\ \dd\theta=\frac{\dd z}{\ii z}$.故原实变函数积分转换为单位圆内的回路积分
    $$
    \begin{aligned}
        I&=\oint_{\abs{z}=1}{\frac{\dd z}{\ii z}\frac{1}{2+\frac{z+z^{-1}}{2}}}\\
        &=\frac{2}{\ii}\oint_{\abs{z}=1}\frac{\dd z}{z^2+4z+1}
    \end{aligned}
    $$
    被积函数$f\left(z\right)$存在两个单极点$z_0=-2\pm\sqrt{3}$，其中$z_0=-2+\sqrt{3}$在单位圆内，其留数为
    $$\Res[f\left(z\right),z]=\lim_{z\to -2+\sqrt{3}}\frac{1}{z-\left(-2-\sqrt{3}\right)}=\frac{1}{2\sqrt{3}}$$
    由留数定理，原积分的值为
    $$I=2\pi\ii\cdot\frac{2}{\ii}\Res[f\left(z\right),z]=\frac{2\pi}{\sqrt{3}}$$
\end{solution}
}
\section{解答题あ（本题8分）}
{
    \begin{exercise}
        求$f\left(z\right)=\frac{1}{z^2\left(z-3\right)}+\sin{\frac{1}{z}}$在$z=0$处的Laurent展开.
    \end{exercise}
    \begin{solution}
        $$
        \begin{aligned}
            f\left(z\right)&=-\frac{1}{3z^2}\frac{1}{1-\frac{z}{3}}+\sin{\frac{1}{z}}\\
            &=-\frac{1}{3z^2}\sum_{k=0}^{\infty}{\left(\frac{z}{3}\right)^k}+\sum_{k=0}^{\infty}{\left(-1\right)^k\frac{z^{2k+1}}{\left(2k+1\right)!}}
        \end{aligned}
        $$
    \end{solution}
}
\section{解答题い（每题10分，共30分）}
{
\begin{exercise}
    利用Laplace变换求解以下ODE.
    $$
    \left\{
        \begin{aligned}
            &y^{\prime\prime}\left(x\right)+y^{\prime}\left(x\right)-6y\left(x\right)=5\ee^{-2t}\\
            &y^{\prime}\left(0\right)=0,\ y\left(0\right)=0
        \end{aligned}
    \right.
    $$
\end{exercise}
\begin{solution}
    设$y\left(x\right)$的Laplace变换为$Y\left(p\right)$，对ODE作Laplace变换，得
    $$p^2Y\left(p\right)+pY\left(p\right)-6Y\left(p\right)=\frac{5}{p+2}$$
    解得变换后函数为
    $$Y\left(p\right)=\frac{5}{\left(p^2-4\right)\left(p-3\right)}=\frac{1}{p+3}+\frac{1/4}{p-2}-\frac{5/4}{p+2}$$
    对$Y\left(p\right)$作Laplace逆变换，得ODE的解为
    $$y\left(x\right)=\ee^{-3t}-\frac{5}{4}\ee^{-2t}+\frac{1}{4}\ee^{2t}$$
\end{solution}

\begin{exercise}
    利用Legendre多项式的生成函数，证明如下递推关系
    $$\left(n+1\right)P_{n+1}\left(x\right)-\left(2n+1\right)xP_n\left(x\right)+nP_{n-1}\left(x\right)=0$$
\end{exercise}
\begin{proof}
    Legendre多项式的生成函数为
    $$\frac{1}{\sqrt{1+r^2-2rx}}=\sum_{l=0}^{\infty}{P_l\left(x\right)r^l}$$
    % \hfill$\square$
\end{proof}

\begin{exercise}
    求$\exp\left(-\frac{x^2}{4a^2}\right)$的Fourier变换，并利用此求解下列PDE
    $$
    \left\{
        \begin{aligned}
            &u_t-a^2u_{xx}=0\\
            &u\left(x,0\right)=\varphi\left(x\right)\\
            &\cdots \mbox{すみません、剩下的边值条件忘了是啥来着……}
        \end{aligned}
    \right.
    $$
\end{exercise}
}
\section{解答题う（本题12分）}
{
    \begin{exercise}
    求解以下PDE.
    $$\left\{
    \begin{aligned}
        &u_t-u_{xx}=\sin{\frac{3\pi}{4}x}\\
        &\mbox{すみません、原题边值条件忘了是啥来着……}\\
    \end{aligned}
    \right.
    $$
    \end{exercise}
}